\documentclass[10pt,a4paper]{amsart}
\usepackage[margin=19mm]{geometry}
\usepackage{amsmath,amssymb,mathtools}
\usepackage[scr=rsfs,cal=euler]{mathalfa}
\usepackage{xfrac}
\usepackage[unicode,hidelinks,bookmarks=false]{hyperref}
\newcommand{\ie}{i.e.\ }
\newcommand{\cf}{cf.\ }
\newcommand{\resp}{respectively\ }
\newcommand{\Subsection}{\subsection}
\providecommand{\nobreakdash}{\nobreak\hbox{-}}
\setlength{\emergencystretch}{2em}
\multlinegap0pt
\usepackage{amssymb}
\usepackage{xcolor}
\usepackage{algorithm}\usepackage{algpseudocode}
\newtheorem{proposition}{Proposition}
\newtheorem{corollary}{Corollary}
\title{A Variational Framework for the Complexity of PDE Solutions}
\author{Juan Esteban Suarez Cardona, Holger Boche, Gitta Kutyniok}
\date{Source: arXiv:2510.21290v3 (CC BY 4.0)}
\begin{document}
\maketitle

\noindent\emph{Source and licence.} A Variational Framework for the Complexity of PDE Solutions. Version arXiv:2510.21290v3 (CC BY 4.0), distributed by arXiv under the Creative Commons Attribution 4.0 International licence, http://creativecommons.org/licenses/by/4.0/, as recorded in the arXiv licence metadata for this exact version and shown on the abstract page https://arxiv.org/abs/2510.21290v3. Full licence notice: \texttt{licenses/arxiv-2510.21290-CC-BY-4.0.txt}.\par\smallskip

\noindent\textit{Editorial scope.} Counted unit: the article's corollary cor:comp\_GF (source0.tex lines 655--661). The article's complete original proof of it is delivered with it (source0.tex lines 662--675, one \texttt{proof} environment of 14 lines). No mathematics is written, repaired or re-derived: the statement and the proof are the article's own lines, in the article's own notation, and no retained line is substituted or re-flowed.  Retained uncounted as the source-local context the proof needs: lines 644--650.  Nothing was omitted from any retained span, so the package carries no omission record.  No retained line contains a citation, so the wrapper carries no bibliography and no bibliography entry.  No retained line contains a figure environment, an image-inclusion command or drawing code, so no image is delivered and the package declares no asset dependency.  Editorial formatting only, in addition: an AMS \texttt{amsart} preamble that restates the article's own declarations and macros used by the retained lines, namely the environments \texttt{proposition}, \texttt{corollary} on separate counters, together with the standard packages those lines need.  The retained environments print in this document as Proposition 1, Corollary 1 in order of appearance, whereas the article prints the same items with the numbering of its own full document.  The complete native source of the article as distributed in the arXiv e-print for this exact version is bundled unchanged as \texttt{sources/187709/source0.tex}.
\begin{proposition}\label{prop:comp_cub}
Let $f\in H^k(\Omega)$ be an $H^k$-computable function, then the $k$-th order Sobolev cubature
$$
I^k_n[f]=\|\mathcal{I}_n[f]\|_{H^k(\Omega)}^2,
$$
is computable. Furthermore, if $f$ is an $H^k$-computable function in polynomial-time, then $I^k_n[f]$ is computable in polynomial-time.
\end{proposition}

\begin{corollary}\label{cor:comp_GF}
    Let $\mathcal{L}^n:C^0(\Omega)\rightarrow \mathbb{R}_+$ be a loss of the form
    \begin{equation}
        \mathcal{L}^n[u]:=\|\mathcal{F}[u_n]-f_n\|_{H^k(\Omega)}^2,
    \end{equation}
    with $f_n\in \Pi_{n}(\Omega)$ a computable polynomial, and $\mathcal{F}:\Pi_{n}(\Omega)\rightarrow \Pi_{n}(\Omega)$ a differentiable operator with $\mathcal{F}[u_n]$ and $\nabla\mathcal{F}[u_n]$ being $C^0$-computable functions. Then, the $k$-th iteration of the explicit Euler gradient flow is $C^0$-computable.  
\end{corollary}

\vspace{-1em}\begin{proof}
    The full gradient of the loss is given by
    \begin{equation}
    \begin{aligned}
        \nabla\mathcal{L}^n[u] &= 2\sum\limits_{|\beta|\leq k}\nabla\mathfrak{F}(u)^T(\mathbb{D^\beta})^T\mathbb{W}\mathbb{D^\beta}(\mathfrak{F}(u)-f_n)\\&
        = 2\nabla\mathfrak{F}(u)^T\mathbb{W}_k(\mathfrak{F}(u)-f_n),
    \end{aligned}
    \end{equation}
    where $\nabla\mathfrak{F}(u):=\nabla\mathcal{F}[u](P_{n}(\Omega))$ and $\mathfrak{F}(u):=\mathcal{F}[u](P_{n}(\Omega))$. By the assumptions on $\mathcal{F}$, both quantities are computable. The result then follows from Proposition~\ref{prop:comp_cub}, noting that since the Sobolev cubature matrices are computable, so is the operator $\mathbb{W}_k = \sum\limits_{|\beta|\leq k}(\mathbb{D}^\beta)^\top \mathbb{W} \mathbb{D}^\beta$. And hence, the iterates 
    $$
   u^{j+1} =u^j-\delta\tau \nabla\mathcal{L}^n[u^{j}],
    $$
    with a computable $\delta\tau>0$, are computable. 
\end{proof}

\end{document}
